\begin{abstract}
\noindent Are pioneers the ones with arrows in their backs? The US
trademark record can answer at census scale: a trademark application is
a dated description of a product a firm expected to sell, and five years
into a registration the owner must swear the product is still in
commerce or lose the mark. Each of 13.99 million applications is scored
on how far its goods/services language sat toward where its industry's
language was moving (its lead) and on how unusual it was for its
industry (its atypicality), and each mark and its owner is followed
through registration, the five-year proof of continued use, the year-ten
renewal, a priced private round, SEC reporting and listing. Across
measures of both survival and the likelihood of reaching financial
markets, being early costs. Among 3.10 million registrations, the most
leading fifth fail the five-year proof $1.4$ points more often than the
most lagging fifth, within Nice class and registration year with
drafting and counsel held ($t = 8.3$), in three of every four classes
and under every partition of the vocabulary tried. The penalty was two
to three times larger for late-1990s filings and returned from 2008; the
exception is 2000--2004, when it reversed inside the four technology
classes: marks still written in the fading dot-com services vocabulary
--- by then the lagging tail --- failed at the era's worst rates, while
leading language was rewarded where the internet met goods, strongest in
e-commerce retail. Unusual language behaves differently. Where a theme
is new only to its Nice class, arriving from elsewhere in the corpus,
its filings are rewarded at the proof; themes new to the entire corpus
are too rare in the record to judge; and the apparent protection of
atypicality is mostly this extension effect. Registration turns on
atypicality --- examination rewards a middling distance from the class's
usual language, and lead does not yet exist to be read --- while the
five-year proof turns on lead. No later rung rewards being early: the
debut's language does not predict a priced round, SEC reporting and
listing fall in lead, and the highest-revenue firms the record can trace
to SEC filings debuted in language their industries already used.

\medskip
\noindent\emph{Keywords:} pioneering; first-mover disadvantage; optimal
distinctiveness; trademarks; product survival; text as data.\\
\emph{JEL:} O31, M13, L10, O34.
\end{abstract}
